% Обзор литературы 0003.
% Сборка: tectonic review.tex (или xelatex + bibtex)
% PDF кладётся в out/, в репозиторий не коммитится.
\documentclass[11pt,a4paper]{article}

\usepackage{iftex}
\ifPDFTeX
  \usepackage[T2A]{fontenc}
  \usepackage[utf8]{inputenc}
  \usepackage[russian]{babel}
\else
  \usepackage{fontspec}
  \usepackage[russian]{babel}
  \IfFontExistsTF{Times New Roman}
    {\setmainfont{Times New Roman}}
    {\IfFontExistsTF{DejaVu Serif}
      {\setmainfont{DejaVu Serif}}
      {\IfFontExistsTF{Liberation Serif}
        {\setmainfont{Liberation Serif}}
        {}}}
\fi

\usepackage[margin=2.4cm]{geometry}
\usepackage{amsmath}
\usepackage{amssymb}
\usepackage{booktabs}
\usepackage{longtable}
\usepackage{array}
\usepackage{enumitem}
\usepackage{hyperref}
\usepackage{microtype}
\setlength{\emergencystretch}{4em}
\sloppy

\hypersetup{
  colorlinks=true,
  linkcolor=black,
  citecolor=black,
  urlcolor=blue
}

\setlist[itemize]{leftmargin=1.4em,itemsep=0.2em,topsep=0.2em}
\setlist[enumerate]{leftmargin=1.6em,itemsep=0.2em,topsep=0.2em}

\title{В каком формате хранить траектории,\\
чтобы обучить предсказатель стоимости решения задачи:\\
обзор литературы}
\author{Обзор 0003}
\date{8 ноября 2026}

\begin{document}
\maketitle

\begin{abstract}
\noindent
Обзор отвечает на вопрос, в каком формате хранить траектории прогонов
LLM-агента, чтобы по ним можно было восстановить обучающую выборку для
предсказателя стоимости решения задачи. Включено 50 работ: схемы траекторий
(ATIF, OpenTelemetry GenAI, AgentOps, структурированные логи), корпуса
траекторий кодовых агентов, object-centric event logs, колоночные форматы,
метаданные и провенанс, а также работы о стоимости, метках и утечке.
Вывод: вопрос решается не одним форматом, а двумя слоями. Первый слой —
неизменяемый поток событий в JSONL, по одному событию на строку, со схемой
уровня ATIF: шаг, вызов LLM, вызов инструмента, наблюдение, границы
контекста и субагенты. Второй слой — производная колоночная таблица
(Parquet) с одной строкой на вызов LLM, в которой собраны признаки для
обучения и сырые счётчики токенов. Обязательны раздельные
\texttt{input}/\texttt{output}/\texttt{reasoning}/\texttt{cached} токены,
модель и версия прайса, момент доступности записи, версия схемы и датасета,
идентификаторы задачи, прогона и родительского шага. Обучающая выборка
строится по cutoff: сначала отсекается будущее, потом считаются признаки;
итоговая стоимость берётся из полного прогона, а незавершённые прогоны
цензурируются, а не помечаются провалом. Готовой работы, которая проверяет
такой формат на downstream-качестве предсказателя стоимости кодового
агента, нет.
\end{abstract}

\section{Вопрос и зачем}

\subsection{Вопрос}

В каком формате хранить траектории прогонов LLM-агента, чтобы по ним можно
было восстановить обучающую выборку для предсказателя стоимости решения
задачи — без утечки будущих признаков, с устойчивой схемой и приемлемой
ценой хранения?

Вопрос распадается на четыре:

\begin{enumerate}
  \item \textbf{Что считать единицей записи.} Прогон, шаг, вызов LLM,
        вызов инструмента или отдельное событие?
  \item \textbf{Какие поля обязательны.} Что нужно, чтобы восстановить и
        стоимость, и признаки, доступные до момента прогноза?
  \item \textbf{Как хранить физически.} Документ на прогон, строка на шаг,
        колоночная таблица или событийный лог?
  \item \textbf{Как из записи получить выборку.} Где проходит cutoff, как
        считаются метки и как не допустить утечки?
\end{enumerate}

Под стоимостью здесь понимается бюджет харнесса: выходные токены, причём
10 входных токенов равны одному выходному, а кеш провайдера не учитывается.
Поэтому формату нужны сырые счётчики по каждому вызову и модель, а не только
одно итоговое число.

\subsection{Зачем}

Обзор 0001 показал, как оценить стоимость задачи до исполнения, а обзор
0002 — как заметить расхождение оценки с фактом и обновить её. Оба обзора
упираются в один и тот же инженерный пробел: чтобы обновлять оценщик по
итогам прогонов, нужно сначала сохранить эти прогоны в пригодном виде.
Сейчас в решении нет структурированной траектории: пишутся итоговые
\texttt{cost} и \texttt{spent} по задаче и хвосты stdout/stderr
(\texttt{specification/behavior/06\_trajectories.md}). Этого достаточно,
чтобы посчитать баллы, но недостаточно, чтобы обучить модель.

Ответ определяет:

\begin{itemize}
  \item схему записи и момент её фиксации — до реализации изменения
        поведения в \texttt{06\_trajectories.md};
  \item какие поля писать всегда, а какие — только в режиме сбора датасета;
  \item нужен ли отдельный формат хранения и агрегат для обучения;
  \item как проверять отсутствие утечки и воспроизводимость выборки.
\end{itemize}

\subsection{Что значит «формат»}

Слово «формат» в литературе означает четыре разных решения, и их полезно
не смешивать:

\begin{enumerate}
  \item \textbf{Единица и схема.} Что является атомом записи и какие поля у
        него есть. Примеры: шаг ATIF~\cite{atif2026}, span
        OpenTelemetry~\cite{otel2026genai}, событие object-centric лога
        OCEL~\cite{berti2024ocel}.
  \item \textbf{Физический контейнер.} JSONL, один JSON-документ, Parquet,
        SQLite, XML.
  \item \textbf{Производные представления.} Полный транскрипт, UI-вид,
        таблица шагов, граф провенанса.
  \item \textbf{Метаданные и версии.} Версия схемы, версия датасета,
        коммит кода, версия харнесса и модели, прайс-лист.
\end{enumerate}

Работы, которые сравнивают «JSON против Parquet», обычно отвечают только на
второй вопрос. Для обучения предсказателя важнее первый и четвёртый: если
шаг не разделяет вызов LLM и наблюдение или не хранит момент доступности
признака, никакой контейнер это не исправит.

\section{Метод}

\subsection{Базы и запросы}

Поиск выполнен 8 ноября 2026 года в arXiv API (поле \texttt{abs:}) и
OpenAlex (поиск по заголовку, аннотации и полному тексту). Просматривались
верхние 1--30 записей на запрос в arXiv и верхние 25 в OpenAlex, сортировка
по релевантности. Язык — английский, годы — 2004--2026; более ранние работы
включались только как метод, на который прямо опираются современные. Точные
запросы, даты и числа — в \texttt{protocol.md}.

Semantic Scholar и Google Scholar не использовались: в обзорах 0001--0002
они вернули HTTP~403. Это ограничивает полноту, особенно индустриальными
отчётами.

\subsection{Критерии отбора}

Включались работы, которые отвечают хотя бы на один из четырёх вопросов:

\begin{itemize}
  \item дают схему или формат представления траектории агента;
  \item дают модель событийных данных с несколькими сущностями и связями;
  \item дают правило построения обучающей выборки для прогноза стоимости
        или предотвращения утечки;
  \item дают механизм провенанса, версионирования или целостности.
\end{itemize}

Допускались LLM- и кодовые агенты, агентные датасеты, event logs и process
mining, наблюдаемость и телеметрия, инженерия данных для ML. Классические
работы включались, если их формат или правило переносится на траекторию
агента. Исключались работы, где траектория упоминается только как
иллюстрация, где речь о траекториях роботов или физических объектов, где
нет схемы или количественной оценки и где формат требует закрытого
проприетарного хранилища.

\subsection{Отбор}

Просмотрено 1\,036 записей: 675 верхних в arXiv, 325 верхних в OpenAlex и
36 по точным названиям. Всего базы вернули 7\,936 совпадений в arXiv и
около 1,8 млн полнотекстовых совпадений в OpenAlex; полная выдача не
просматривалась, поэтому в отбор вошли только верхние по релевантности
записи. После удаления 33 дубликатов просмотрено 1\,003 заголовка и
аннотации. На скрининге исключено 933 записи, для детального разбора
отобрано 70, после разбора исключено 20. В обзор включено 50 работ. Схема
и причины исключения — в \texttt{sources.md}.

\subsection{Извлечение и оценка качества}

Для каждой работы выписаны: единица хранения, схема и обязательные поля,
физический формат, работа с вложенностью и связями, метки и стоимость,
провенанс и версии, количественный результат, ограничение и вывод для нас.
Надёжность оценивалась по рецензируемости, наличию публичной спецификации,
воспроизводимости, размеру данных, сравнению с baseline и явной обработке
утечки. Препринты, стандарты и технические отчёты помечены отдельно.
Мета-анализ не проводился: форматы и метрики несопоставимы.

\section{Что известно}

\subsection{Единица записи: прогон, шаг или вызов}

Все корпуса траекторий сходятся в том, что минимальная полезная единица —
не прогон, а \emph{шаг}: одно решение агента вместе с его действиями и
наблюдениями. ATIF определяет шаг как system-, user- или agent-ход, внутри
которого есть сообщение, рассуждение, один или несколько вызовов
инструментов и наблюдение~\cite{atif2026}. Исследование
Thought-Action-Result сводит логи трёх SE-агентов к той же
тройке~\cite{bouzenia2025tar}; OpenDiscoveryTrace пишет на шаг девять полей,
включая ошибки и триггеры ревизии~\cite{bansal2026odt}; TraceLab считает
шаги и вызовы инструментов раздельно~\cite{zhu2026tracelab}.

Но для прогноза стоимости шага недостаточно. TokenCast показывает, что
стоимость сегмента складывается из двух величин: собственного расхода и
прироста контекста, который этот сегмент вносит в последующие
вызовы~\cite{ouyang2026tokencast}. Если хранить только шаг, внутри которого
было несколько вызовов LLM, прирост контекста между ними теряется. Поэтому
единица записи должна быть иерархической:

\begin{center}
прогон $\rightarrow$ шаг $\rightarrow$ вызов LLM $\rightarrow$ вызов
инструмента $\rightarrow$ наблюдение.
\end{center}

ATIF решает это полем \texttt{llm\_call\_count}: если вызовов несколько,
метрики шага агрегированы и per-call атрибуция недоступна; если ноль —
шаг детерминированный и метрик быть не должно~\cite{atif2026}. Это честное
признание компромисса, но для обучения предсказателя стоимости агрегат
шага недостаточен: нужна отдельная запись на каждый вызов LLM.

\subsection{Схема: какие поля обязательны}

Литература даёт три уровня полей.

\paragraph{Слой действия.} Кто и что сделал: \texttt{source},
\texttt{message}, \texttt{reasoning\_content}, \texttt{tool\_calls},
\texttt{observation}, ссылки на родительский шаг и на вызов, породивший
наблюдение~\cite{atif2026,bouzenia2025tar}. AgentTrace добавляет
разделение operational, cognitive и contextual поверхностей: что агент
сделал, что он «думал» и в каком контексте~\cite{alsayyad2026agenttrace}.
Agent Execution Record идёт дальше и делает первого класса намерение,
наблюдение, вывод, версию плана, цепочку свидетельств и уверенность
вердикта~\cite{vispute2026aer}.

\paragraph{Слой стоимости.} Сколько стоил каждый вызов: модель, провайдер,
входные, выходные, reasoning- и cached-токены, время начала и конца,
статус, ошибка. ATIF хранит \texttt{prompt\_tokens},
\texttt{completion\_tokens}, \texttt{cached\_tokens}, \texttt{cost\_usd},
token ids и logprobs~\cite{atif2026}. OpenTelemetry GenAI стандартизует
\texttt{gen\_ai.usage.input\_tokens} и
\texttt{gen\_ai.usage.output\_tokens}, причём входные включают cached, а
выходные — reasoning; контент можно хранить отдельно, а в span держать
ссылку \texttt{gen\_ai.input.messages\_ref}~\cite{otel2026genai}. Это
важное разделение: счётчики должны быть в записи всегда, а тяжёлый контент
может лежать рядом.

\paragraph{Слой контекста.} Что было в контексте до вызова и как он вырос.
ATIF вводит \texttt{context\_management} с типами \texttt{compaction},
\texttt{pruning}, \texttt{injection} и границей \texttt{replace},
\texttt{append}, \texttt{truncate}; после границы \texttt{replace}
предыдущие шаги сохраняются для аудита, но не входят в контекст
модели~\cite{atif2026}. Это ровно та информация, без которой нельзя
пересчитать стоимость и нельзя понять, какие признаки были доступны.

Проверка полноты схемы — отдельная проблема. Solozobov применяет один и
тот же Decision Event Schema к шести вендорским режимам и находит, что
полнота strict-governance свойств колеблется от 42,9\% до 85,7\%, а
reasoning trace не заполняется ни в одном режиме~\cite{solozobov2026reconstruct}.
TelemetrySuffBench показывает ещё резче: при полной телеметрии origin-step
Top-1 составляет 33,8--97,2\%, а metadata-, OpenTelemetry- и
OpenInference-совместимые виды сохраняют 99,5--100\% detection F1, но не
более 0,5\% точности локализации причины; удаление decision content
обнуляет локализацию~\cite{zhu2026telemetrysuff}. Отсюда правило: формат
должен явно помечать отсутствующие поля и не притворяться, что
«структурированный лог» уже содержит решение.

\subsection{Физический формат: документ, таблица или событийный лог}

\paragraph{JSONL как канонический слой.}
VCC называет траекторию структурированным документом: сессия кодового
агента содержит пользовательские ходы, текст ассистента, chain of thought,
вызовы инструментов, результаты, вызовы субагентов, границы компакции и
директивы харнесса, а размер может превышать десять тысяч строк JSONL.
VCC хранит lossless-транскрипт и строит поверх него UI- и adaptive-виды;
на AppWorld это даёт +1,1--4,2 п.п. pass rate и снижает токены рефлектора
на половину-две трети~\cite{zhang2026vcc}. Практический вывод: сырой JSONL
надо сохранять целиком, а производные представления строить как проекции,
а не заменять им исходник.

\paragraph{Parquet как слой обучения.}
Колоночные форматы выигрывают там, где нужны сканы и соединения.
Сравнение Parquet и ORC показывает, что dictionary encoding по умолчанию,
скорость декодирования важнее степени сжатия, а блочное сжатие стоит
делать опциональным; при этом форматы имеют известные неэффективности на
ML-нагрузках~\cite{zeng2023columnar}. ParquetDB демонстрирует, что Parquet
держит complex и nested типы, поддерживает predicate pushdown и обходит
SQLite и MongoDB на больших объёмах; проверка — 4,8 млн вложенных
записей~\cite{lang2025parquetdb}. Lance добавляет важный нюанс: Parquet при
правильной настройке даёт более чем 60-кратный рост random access, но за
счёт scan и RAM; альтернативная структурная кодировка убирает этот
компромисс~\cite{pace2025lance}. Параллельная запись вложенных данных
переменного размера возможна, но требует формата с явными offsets, а не
склейки строк~\cite{hahnfeld2024nested}.

\paragraph{Эпизодный стандарт.}
RLDS задаёт lossless-формат эпизода, не зависящий от исходного формата, с
записью, replay, аннотированием и шарингом через TFDS~\cite{ramos2021rlds}.
Для траекторий агента это хороший ориентир: эпизод — единица, шаги —
вложенные записи, формат не теряет информацию при конвертации.

\paragraph{Event log как слой связей.}
XES задаёт минимум: case id, activity, timestamp и атрибуты, а расширения
позволяют добавлять новые понятия без потери
семантики~\cite{wynn2024xes}. JXES показывает, что JSON-представление
возможно, но производительность зависит от парсера~\cite{shankara2020jxes}.
OCEL 2.0 идёт дальше: событие ссылается на несколько объектов, у объектов
есть жизненный цикл, связи и квалифицированные отношения; обмен — SQLite,
XML или JSON~\cite{berti2024ocel}. Extended Event Log пытается объединить
признаки нескольких стандартов~\cite{suleiman2024extended}. Для траектории
агента это ближе, чем XES: задача, сессия, шаг, вызов инструмента и
субагент — разные объекты, и связь между ними многие-ко-многим.

\paragraph{Что выбрать.}
Ни один контейнер не решает все задачи. JSONL хорош для lossless-потока,
аудита и replay, но плох для сканов и соединений. Parquet хорош для
обучения, но теряет порядок и сырой контент, если писать только его.
OCEL/XES хороши для связей, но тяжелы, а XES вообще один кейс на событие и
не поддерживает вложенных сущностей. OTel spans хороши для межвендорского
учёта, но это не lossless-транскрипт: контент часто выносится в
\texttt{\_ref}, а сэмплирование и объём ограничены. Отсюда гибрид:
неизменяемый JSONL-поток как источник истины и Parquet-таблица как
производный слой обучения.

\subsection{Многосущностные связи}

Простая строка «прогон — шаг — токены» теряет структуру, которая нужна и
для анализа, и для обучения. Обзор execution provenance предлагает
типизированный граф прогона: узлы — шаги, вызовы, артефакты, свидетельства;
рёбра — причинность, поддержка, делегирование; шесть измерений таксономии
включают гранулярность, время, форму представления и trust-функции. Авторы
называют три пробела: нет интеграции тонкого провенанса с trust-функциями,
нет бенчмарка полной цепочки и мало tamper-evident трасс~\cite{wang2026trust}.

ATIF отвечает на это ссылками: \texttt{tool\_call\_id} связывает вызов с
наблюдением, \texttt{subagent\_trajectories} и
\texttt{trajectory\_id} — родителя с подагентом,
\texttt{continued\_trajectory\_ref} — сегменты одного прогона, а
\texttt{is\_copied\_context} помечает шаги, скопированные через границу
компакции~\cite{atif2026}. OCEL даёт более общий язык для тех же связей:
объекты, события, отношения и их изменения во времени~\cite{berti2024ocel}.
Для хранения это означает, что нужны как минимум три идентификатора:
\texttt{task\_id}, \texttt{run\_id}, \texttt{step\_id}, плюс
\texttt{call\_id} и \texttt{parent\_step\_id}. Без них нельзя ни собрать
историю, ни исключить копированный контекст, ни связать расход с
конкретным действием.

Унификация формата — отдельный инженерный слой. AgentOhana агрегирует
разнородные multi-turn траектории и приводит их к согласованному
формату с общим data loader~\cite{zhang2024agentohana}; ActionStudio
делает то же через Unified Format 2.0 и сообщает до 9-кратного роста
throughput обучения~\cite{zhang2025actionstudio}. Это аргумент против
ad-hoc JSON: если формат не стандартизован, каждая новая среда добавляет
свой конвертер и свою потерю информации.

\subsection{От траектории к обучающей выборке}

\paragraph{Признаки должны быть доступны на cutoff.}
LEAP формализует ранний прогноз как задачу с ограничением доступности:
сначала отсечь будущее по cutoff, потом делать joins и агрегации, а не
наоборот; провенанс признаков проверяется отдельно. На OULAD без этого
ранние метрики завышены, а заметный прирост появляется только после
третьей недели наблюдения~\cite{le2026leap}. Для предсказателя стоимости
это означает: строка обучающей выборки — это не «все поля прогона», а
снимок на момент $t$, где $t$ — начало решения или начало шага.

\paragraph{Момент доступности важнее timestamp.}
Аудит point-in-time утечки в энергетическом рынке показывает, что выбор
«последнего прогона до гейта» в более чем 99\% дней для 6 из 10
конфигураций берёт запись, которая ещё не была опубликована; у GFS доля
меняется с 0,6\% зимой до 99,8\% летом, а ревизии задним числом вообще не
сохраняются~\cite{alfalah2026lookahead}. Аналогия прямая: у траектории
недостаточно хранить \texttt{timestamp} события. Нужен
\texttt{available\_at} — момент, когда запись стала видна сборщику, — и
история ревизий, иначе выборка будет использовать будущее.

\paragraph{Признаки материализуются с задержкой.}
Feature Freshness Budgets вводит freshness gap между событием и
видимостью признака в online feature store и формальный per-feature budget
со staleness bound и closed-form порогом; предсказание подтверждается на
Kafka, Redis и PostgreSQL~\cite{rajula2026freshness}. Если признаки
считаются заранее, в записи нужны версия признака и задержка его
материализации; training-serving skew — это свойство данных, а не только
модели.

\paragraph{Метка: цена полного прогона, а не шага.}
TokenCast предсказывает остаток по ходу, но обучается на сегментах,
у которых есть собственный расход и прирост контекста~\cite{ouyang2026tokencast}.
CRATER показывает, что незавершённый при лимите ответ — это цензурированное
наблюдение, а не «неправильный»: корректное и некорректное завершение —
конкурирующие события, и вероятность каждого считается отдельно~\cite{jiang2026crater}.
Для нашей задачи это значит: \texttt{status} должен различать
\texttt{solved}, \texttt{wrong}, \texttt{unfinished},
\texttt{budget\_exhausted}, \texttt{error}; \texttt{unfinished} нельзя
сваливать в \texttt{wrong}.

\paragraph{Утечка — свойство пайплайна.}
LeakageDetector ищет утечку в коде ML-пайплайна и предлагает
исправление~\cite{alomar2025leakagedetector}. Проверку стоит встроить в
сборку датасета. Дополнительные правила для траекторий: не смешивать
прогоны одной задачи в train и test; не использовать поля, появившиеся
после cutoff; исключать шаги с \texttt{is\_copied\_context}; не считать
признаком итоговый \texttt{cost}, если он не был известен на момент
прогноза.

\paragraph{Разметка шага.}
Annotrace показывает, что шаг — правильная единица разметки, а форматы
OTel GenAI, LangSmith и Langfuse приводятся к общему вложенному списку
шагов; на корпусе multi-agent QA инструмент воспроизводит 235 единиц
ручной разметки, а две модели-аннотатора совпадают в 91\% при
$\kappa=0{,}82$~\cite{huang2026annotrace}. Если метка «сколько стоил шаг»
или «был ли шаг полезен» ставится вручную, её надо хранить на том же шаге,
а не в отдельном CSV без ключей.

\subsection{Стоимость: сырые счётчики, а не одно число}

Три работы объясняют, почему итогового \texttt{cost} недостаточно.

Во-первых, стоимость агентной задачи в 1000 раз больше, чем code reasoning
или code chat, разброс одного и того же задания достигает 30 раз, входные
токены доминируют, а самооценка модели имеет корреляцию с фактом не выше
0,39 и систематически занижена~\cite{bai2026money}. Значит, нужен
распределительный прогноз и шаговые признаки, а не точечная оценка.

Во-вторых, расход распределён по фазам неравномерно: в ChatDev на 30
задачах Code Review даёт 59,4\% токенов, а input — 53,9\%~\cite{salim2026tokenomics}.
Метка фазы или активности на шаге — дешёвый и полезный признак.

В-третьих, стоимость надо атрибутировать шагу. ClawTrace записывает
cost-attributed traces и TraceCard; на SpreadsheetBench и SkillsBench
удаление prune-патчей втрое увеличивает регрессии качества, не снижая
медианную стоимость, а агрегатная стоимость вообще не меняется~\cite{yuan2026clawtrace}.
Без per-step атрибуции нельзя отличить полезный шаг от дорогого
бесполезного.

Отсюда практическое правило: хранить сырые счётчики и модель, а стоимость
считать при чтении по версионируемому прайсу. Two-Ledger разделяет
settlement-журнал (FOCUS) и attribution-журнал invocation-grain на OTel
GenAI, сшивает их по trace context, A2A/MCP, бизнес-ключу и вероятностному
совпадению, применяет versioned rate card и получает покрытие выше 99\%
при остатке 0,24 п.п.~\cite{acharya2026twoledger}. ATIF сам предупреждает,
что \texttt{cost\_usd} — снимок на момент исполнения, потому что цены
меняются~\cite{atif2026}; для бюджета соревнования с весами 10:1 и без
кеша это особенно важно — пересчёт должен быть возможен из сырых чисел.

И наконец, «меньше токенов» не значит «меньше денег». Сжатие tool-output
на 38,4\% увеличило billed cost на 6,8\%, корреляция token reduction и
cost reduction составила $r=0,15$, а cache creation/reads доминировали во
входной стоимости~\cite{weinberger2026tokenreduction}. Поэтому в записи
нужны \texttt{cached\_tokens}, \texttt{cache\_creation\_tokens} и
фактический биллинг, если он доступен.

\subsection{Провенанс, версии и целостность}

\paragraph{Трассу нельзя оставлять только агенту.}
Qin и соавторы показывают, что Claude Code, Codex, Antigravity, Open Code
и Grok Build позволяют агенту удалить свои трассы по просьбе; исключение —
Muse Code. Внешний атакующий может это использовать, а поведение возникает
при оптимизации reward~\cite{qin2026tamper}. Вывод для формата: писать
траекторию должен внешний append-only приёмник, а не только процесс
агента; в идеале — с хэш-цепочкой.

\paragraph{Хэш-цепочка и content-addressed схемы.}
Agent-Native Telemetry строит State-Delta Evidence Ledger из четырёх
примитивов (Transitions, Observations, Relations, State Checkpoints),
подписывает и хэширует батчи, а на AIOpsLab и OTel Astronomy Shop снижает
payload на 96,4\%, контекст LLM на 88,8\% и число запросов на 66,2\%;
500 из 500 подмен хранилища обнаружены~\cite{he2026atp}. Traccia добавляет
хэшированный trace ledger и SHA-256-отпечатки, привязанные к требованиям
EU AI Act~\cite{naik2026traccia}. Это не обязательно для обучения, но
дёшево и резко повышает доверие к датасету.

\paragraph{Версии данных и метаданные.}
Dataset Management Platform делает storage engine источником истины с
версиями, доступом и снапшотами для обучения и оценки~\cite{mao2023datasetmgmt}.
TableVault добавляет конкурентное исполнение, воспроизводимость и
версионирование артефактов в LLM-пайплайнах~\cite{zhao2025tablevault}.
Croissant задаёт единое описание датасета, которое понимают инструменты и
платформы, и уже поддерживается репозиториями с сотнями тысяч
датасетов~\cite{akhtar2024croissant}. MLflow2PROV извлекает
W3C-PROV-совместимый граф из Git и MLflow~\cite{schlegel2023mlflow2prov}.
Аудит почти 4\,000 датасетов показывает, что цепочки деривации и лицензии
почти не прослеживаются: менее 33\% датасетов имеют ограничительную
лицензию, но более 80\% исходного контента — non-commercial
ограничения~\cite{longpre2024provenance}. Для траекторий это означает:
поля \texttt{source}, \texttt{license}, \texttt{code\_commit},
\texttt{harness\_version}, \texttt{model\_version}, \texttt{schema\_version},
\texttt{dataset\_version} должны быть обязательными метаданными, а не
комментарием в README.

\subsection{Классические аналоги}

Event logs и process mining дают два урока, которых нет в LLM-литературе.
Первый: тип лога меняет постановку прогноза. Сравнительный обзор
predictive process monitoring показывает, что методы для традиционных и
object-centric логов различаются, и object-centric лог даёт больше
признаков, но требует другой агрегации~\cite{fioretto2025ppm}. Второй:
качество лога — самостоятельная проблема. Опрос практиков выявляет самые
частые и важные проблемы и разрыв между воспринимаемой частотой и
важностью~\cite{dakic2023quality}; систематический обзор генерации логов
находит 29 проблем и 5 категорий требований, включая data engineering и
domain knowledge engineering~\cite{elalaoui2026eventlog}. Это ровно то,
что обычно пропускают при сборе траекторий: лог считается побочным
продуктом, а не данными с требованиями к качеству.

\section{Сравнение}

Таблица~\ref{tab:sources} сводит включённые работы. Строки сгруппированы
по вкладу: схема траектории, корпус, event log, физическое хранение,
метаданные, стоимость и утечка.

{\scriptsize
\setlength{\tabcolsep}{3pt}
\renewcommand{\arraystretch}{1.12}
\begin{longtable}{@{}>{\raggedright\arraybackslash}p{0.55cm}>{\raggedright\arraybackslash}p{3.0cm}>{\raggedright\arraybackslash}p{2.7cm}>{\raggedright\arraybackslash}p{4.3cm}>{\raggedright\arraybackslash}p{3.0cm}@{}}
\caption{Включённые работы.}\label{tab:sources}\\
\toprule
ID & Работа & Что даёт & Данные и результат & Ограничение \\
\midrule
\endfirsthead
\toprule
ID & Работа & Что даёт & Данные и результат & Ограничение \\
\midrule
\endhead
\bottomrule
\endfoot

S01 & Li~\cite{atif2026} & схема шага и метрик & JSON; per-step tokens/cost, \texttt{llm\_call\_count}, \texttt{is\_copied\_context}, субагенты & спецификация, нет downstream \\
S02 & Dong et al.~\cite{dong2024agentops} & таксономия наблюдаемости & артефакты и данные жизненного цикла & нет схемы полей \\
S03 & AlSayyad et al.~\cite{alsayyad2026agenttrace} & слои логов & operational, cognitive, contextual & нет замера \\
S04 & Zhang, Agrawala~\cite{zhang2026vcc} & lossless-транскрипт + виды & JSONL; +1,1--4,2 п.п. pass, токены рефлектора $-50$--$67\%$ & один продукт \\
S05 & He, Yu~\cite{he2026atp} & state-delta телеметрия & 4 примитива, hash-chain; $-96{,}4\%$ payload, $-88{,}8\%$ контекста & инфраструктура \\
S06 & Vispute, Kadam~\cite{vispute2026aer} & reasoning provenance & intent, observation, inference, планы, evidence chains & нет эмпирики \\
S07 & Wang et al.~\cite{wang2026trust} & граф провенанса & 6 измерений, 9 отношений, 3 пробела & обзор \\
S08 & Zhu, Pu~\cite{zhu2026telemetrysuff} & достаточность телеметрии & detection F1 99,5--100\%, origin Top-1 $\le 0{,}5\%$ для OTel-видов & синтетика \\
S09 & Solozobov~\cite{solozobov2026reconstruct} & полнота свойств & 42,9--85,7\% strict-governance; reasoning не заполняется & один аннотатор \\
S10 & Huang et al.~\cite{huang2026annotrace} & разметка шага & 6 форматов; 235 единиц; согласие 91\%, $\kappa=0{,}82$ & инструмент \\
S11 & Pilditch et al.~\cite{pilditch2026transect} & длинные прогоны & таймлайн, ~13 млн токенов, экспорт таблиц & оценка \\
S12 & OpenTelemetry~\cite{otel2026genai} & межвендорский словарь & usage tokens, conversation id, messages\_ref, tool definitions & Development \\
S13 & Manam et al.~\cite{manam2026otel} & связка ops и eval & 8\,500 spans/s, 12\,000 metrics/s; F1 0,84 & эталонная конфигурация \\
S14 & Naik et al.~\cite{naik2026traccia} & governance-журнал & OTel + SHA-256, EU AI Act & нет замера \\
S15 & Ahmad et al.~\cite{ahmad2026openswe} & корпус SE & 207\,489 траекторий, 9 языков; 61,7\% SWE-bench Verified & нет cost-схемы \\
S16 & Bansal, Balaji~\cite{bansal2026odt} & 9 полей на шаг & 558 траекторий, 124 задачи, 7 моделей & не стоимость \\
S17 & Baumann et al.~\cite{baumann2026swechat} & реальные сессии & 18 тыс. сессий, 2 млн вызовов; 59\% кода доходит до коммита & нет стоимости \\
S18 & Zhu et al.~\cite{zhu2026tracelab} & serving-трасса & 4\,300 сессий, 350 тыс. шагов, 430 тыс. вызовов & не обучение \\
S19 & Bouzenia, Pradel~\cite{bouzenia2025tar} & TAR-триплет & 120 траекторий, 2\,822 взаимодействия & малая выборка \\
S20 & Zhang et al.~\cite{zhang2024agentohana} & унификация формата & общий data loader, xLAM-v0.1 & нет схемы в аннотации \\
S21 & Zhang et al.~\cite{zhang2025actionstudio} & Unified Format 2.0 & до 9$\times$ throughput, 98k траекторий & нет сравнения форматов \\
S22 & Berti et al.~\cite{berti2024ocel} & object-centric лог & объекты, события, связи; SQLite/XML/JSON & нет LLM-замера \\
S23 & Wynn et al.~\cite{wynn2024xes} & XES & case id, activity, timestamp, расширения & один кейс на событие \\
S24 & Shankara Narayana et al.~\cite{shankara2020jxes} & JSON для XES & 4 парсера, разные скорости & нет downstream \\
S25 & Suleiman, Kassem~\cite{suleiman2024extended} & единый event log & объединение стандартов & нет эмпирики \\
S26 & Fioretto, Masciari~\cite{fioretto2025ppm} & сравнение PPM & traditional vs object-centric & обзор \\
S27 & Dakić et al.~\cite{dakic2023quality} & качество логов & частые проблемы и предобработка & опрос \\
S28 & El Alaoui Talibi et al.~\cite{elalaoui2026eventlog} & генерация логов & 29 проблем, 5 категорий требований & SLR \\
S29 & Zeng et al.~\cite{zeng2023columnar} & Parquet/ORC & dictionary encoding, скорость декодирования, ML-неэффективности & нет траекторий \\
S30 & Lang et al.~\cite{lang2025parquetdb} & ParquetDB & nested типы, predicate pushdown; 4,8 млн записей & нет обучения \\
S31 & Pace et al.~\cite{pace2025lance} & random access & Parquet $>60\times$ при настройке; Lance без компромисса & нет траекторий \\
S32 & Hahnfeld et al.~\cite{hahnfeld2024nested} & параллельная запись & nested variable-size, ROOT RNTuple & HEP \\
S33 & Ramos et al.~\cite{ramos2021rlds} & эпизодный стандарт & lossless, replay, TFDS & RL, не LLM \\
S34 & Akhtar et al.~\cite{akhtar2024croissant} & метаданные датасета & единое описание, сотни тысяч датасетов & не данные \\
S35 & Qin et al.~\cite{qin2026tamper} & целостность трасс & агенты удаляют свои трассы; исключение Muse Code & локальные harness \\
S36 & Mao et al.~\cite{mao2023datasetmgmt} & версии данных & source of truth, снапшоты & технический отчёт \\
S37 & Zhao, Krishnan~\cite{zhao2025tablevault} & версии артефактов & конкурентность, воспроизводимость & нет агентских траекторий \\
S38 & Schlegel, Sattler~\cite{schlegel2023mlflow2prov} & W3C PROV & граф из Git и MLflow & только MLflow \\
S39 & Longpre et al.~\cite{longpre2024provenance} & провенанс датасетов & ~4\,000 датасетов; $>80\%$ non-commercial & аудит \\
S40 & Ouyang et al.~\cite{ouyang2026tokencast} & прогноз стоимости & сегмент = расход + прирост контекста; MAE $-14{,}5\%$, $-21{,}3\%$ токенов & один агент \\
S41 & Bai et al.~\cite{bai2026money} & расход агентов & 8 моделей, SWE-bench; $1000\times$ токенов, разброс $30\times$, $r\le0{,}39$ & SWE-bench \\
S42 & Salim et al.~\cite{salim2026tokenomics} & фазы SDLC & Code Review 59,4\%, input 53,9\% & 30 задач \\
S43 & Yuan et al.~\cite{yuan2026clawtrace} & per-step cost & TraceCard; prune-патчи как guardrails & узкие бенчмарки \\
S44 & Acharya~\cite{acharya2026twoledger} & два журнала & FOCUS + OTel; coverage $>99\%$, residual 0,24 п.п., p99 $-61\%$ & симуляция \\
S45 & Weinberger, Hozez~\cite{weinberger2026tokenreduction} & токены $\ne$ деньги & $-38{,}4\%$ токенов, $+6{,}8\%$ billed cost, $r=0{,}15$ & Claude Code \\
S46 & Le et al.~\cite{le2026leap} & cutoff против утечки & cutoff-first truncation, аудит провенанса; OULAD & LMS-логи \\
S47 & AlOmar et al.~\cite{alomar2025leakagedetector} & детектор утечки & статический анализ ML-кода & не данные \\
S48 & Alfalah~\cite{alfalah2026lookahead} & point-in-time утечка & 128\,656 прогонов; $>99\%$ неопубликованных записей & энергетика \\
S49 & Rajula~\cite{rajula2026freshness} & freshness budget & staleness bound, Kafka/Redis/Postgres & инфраструктура \\
S50 & Jiang et al.~\cite{jiang2026crater} & цензурирование & correct/wrong/unfinished, ожидаемый расход & reasoning \\
\end{longtable}
}

\section{Пробел и вывод}

\subsection{Пробел}

Литература хорошо отвечает на три вопроса по отдельности:

\begin{itemize}
  \item \emph{что записывать} — шаг, вызов, наблюдение, метрики, контекст;
  \item \emph{как хранить} — JSONL для lossless-потока, Parquet для сканов,
        object-centric лог для связей;
  \item \emph{как не допустить утечки} — cutoff, point-in-time, провенанс,
        цензурирование.
\end{itemize}

Но нет работы, которая бы проверила конкретный формат траектории кодового
агента на downstream-качестве предсказателя стоимости при общем бюджете
токенов. Существующие корпуса (Open-SWE-Traces, OpenDiscoveryTrace,
SWE-chat, TraceLab) собираются для обучения агента или анализа поведения,
а не для прогноза стоимости; в них либо нет per-call счётчиков, либо нет
момента доступности, либо нет цены. Работы по стоимости (TokenCast,
Tokenomics, ClawTrace, Two-Ledger) показывают, какие поля нужны, но не
дают готовой схемы хранения. ATIF ближе всех, но это спецификация, а не
эмпирическая проверка; к тому же она хранит \texttt{cost\_usd} снимком и
допускает агрегацию нескольких вызовов в один шаг.

Пробел узкий и практический: нужен формат, который одновременно

\begin{enumerate}
  \item lossless для аудита и replay;
  \item пригоден для сканов и соединений при обучении;
  \item хранит сырые счётчики, достаточные для пересчёта стоимости по
        правилам харнесса;
  \item позволяет построить cutoff-выборку без утечки;
  \item версионирован и защищён от потери.
\end{enumerate}

\subsection{Рекомендуемый формат}

Ответ на вопрос — два слоя и метаданные, а не один файл.

\paragraph{Слой 1. Канонический поток событий.}
Неизменяемый JSONL, одна строка — одно событие. Схема совместима с ATIF:
прогон, шаг, вызов LLM, вызов инструмента, наблюдение, граница контекста,
субагент. Обязательные поля шага: \texttt{schema\_version},
\texttt{task\_id}, \texttt{run\_id}, \texttt{step\_id},
\texttt{parent\_step\_id}, \texttt{source}, \texttt{timestamp},
\texttt{available\_at}, \texttt{model}, \texttt{provider},
\texttt{input\_tokens}, \texttt{output\_tokens},
\texttt{reasoning\_tokens}, \texttt{cached\_tokens},
\texttt{cache\_creation\_tokens}, \texttt{tool\_name},
\texttt{status}, \texttt{error\_type}. Контент сообщений и наблюдений
хранится в этом же событии, если он не слишком велик; иначе — по ссылке
(как \texttt{messages\_ref} в OTel~\cite{otel2026genai}). Поток пишется
вне песочницы агента, append-only, по возможности с хэш-цепочкой
(как в ATP~\cite{he2026atp} и Traccia~\cite{naik2026traccia}).

Почему JSONL, а не сразу Parquet: VCC показывает ценность lossless-
транскрипта и производных видов~\cite{zhang2026vcc}; RLDS — ценность
lossless-эпизода~\cite{ramos2021rlds}; ATIF — ценность читаемой схемы,
которую можно валидировать Pydantic-моделями~\cite{atif2026}. JSONL также
устойчив к сбою: потерянная строка не ломает весь файл, тогда как
битый Parquet-блок может.

\paragraph{Слой 2. Таблица обучения.}
Parquet, одна строка на вызов LLM (не на шаг). Колонки — плоские:
идентификаторы, модель, времена, счётчики токенов, длина контекста до
вызова, прирост контекста, тип инструмента, задержка, статус, фаза,
цена. Вложенные \texttt{tool\_calls} и \texttt{observation} остаются в
JSONL; в таблицу попадают только те поля, которые нужны модели. Parquet
даёт predicate pushdown и nested типы, если они понадобятся, и хорошо
сканируется~\cite{lang2025parquetdb,zeng2023columnar}. При росте корпуса
до миллионов строк стоит проверить random access: Parquet настраивается,
а Lance — альтернатива~\cite{pace2025lance}.

Разделение слоёв важно: JSONL — источник истины и audit trail, Parquet —
материализованное представление, которое можно пересобрать. Если Parquet
потерян, его можно восстановить; если потерян JSONL, восстановить нечего.

\paragraph{Метаданные.}
Рядом с данными — описание в духе Croissant: поля, типы, единицы,
происхождение, лицензия~\cite{akhtar2024croissant}. Отдельно —
\texttt{dataset\_version}, \texttt{schema\_version}, \texttt{code\_commit},
\texttt{harness\_version}, \texttt{model\_version}, \texttt{rate\_card\_version}.
Версия прайса обязательна: ATIF прямо говорит, что цены меняются, а
Two-Ledger показывает, как считать стоимость по versioned rate
card~\cite{atif2026,acharya2026twoledger}. Для бюджета соревнования
(10 входных = 1 выходной, кеш не учитывается) прайс — это правило
пересчёта, а не рыночная цена.

\paragraph{Обучающее представление.}
Строка выборки — снимок на cutoff $t$: все признаки, у которых
\texttt{available\_at} $\le t$; метка — стоимость остатка или полного
прогона, посчитанная по сырым счётчикам. Правила:

\begin{enumerate}
  \item сначала truncation по cutoff, потом joins и агрегации
        (LEAP~\cite{le2026leap});
  \item \texttt{unfinished} и \texttt{budget\_exhausted} — цензурированные
        исходы, а не \texttt{wrong} (CRATER~\cite{jiang2026crater});
  \item шаги с \texttt{is\_copied\_context} исключаются из обучения
        (ATIF~\cite{atif2026});
  \item split — по задаче или репозиторию, не по шагам, иначе один прогон
        попадает в train и test;
  \item проверка утечки — автоматическая, в пайплайне сборки
        (LeakageDetector~\cite{alomar2025leakagedetector});
  \item версия признака и задержка материализации сохраняются
        (Feature Freshness~\cite{rajula2026freshness}).
\end{enumerate}

\paragraph{Чего не делать.}
Не хранить только итоговый \texttt{cost} и \texttt{spent}: этого мало для
обучения (TokenCast~\cite{ouyang2026tokencast}). Не хранить только шаг с
агрегированными метриками: теряется per-call атрибуция
(ATIF~\cite{atif2026}). Не хранить только Parquet: теряется порядок и
сырой контекст (VCC~\cite{zhang2026vcc}). Не хранить трассу только внутри
песочницы агента: агент может её удалить (Qin et al.~\cite{qin2026tamper}).
Не считать «меньше токенов» синонимом «меньше денег»
(Weinberger, Hozez~\cite{weinberger2026tokenreduction}).

\subsection{Минимальная схема}

Таблица~\ref{tab:fields} перечисляет минимальный набор полей. Она не
заменяет спецификацию изменения, но задаёт проверяемый минимум.

{\small
\begin{longtable}{@{}>{\raggedright\arraybackslash}p{4.0cm}>{\raggedright\arraybackslash}p{4.2cm}>{\raggedright\arraybackslash}p{6.2cm}@{}}
\caption{Минимальные поля записи.}\label{tab:fields}\\
\toprule
Поле & Где & Зачем \\
\midrule
\endfirsthead
\toprule
Поле & Где & Зачем \\
\midrule
\endhead
\bottomrule
\endfoot

\texttt{schema\_version} & прогон & валидация и миграции \\
\texttt{dataset\_version} & датасет & воспроизводимость выборки \\
\texttt{task\_id}, \texttt{run\_id} & прогон & группировка и split \\
\texttt{step\_id}, \texttt{parent\_step\_id} & шаг & порядок и вложенность \\
\texttt{call\_id} & вызов & связь с наблюдением \\
\texttt{source} & шаг & system/user/agent \\
\texttt{timestamp}, \texttt{available\_at} & шаг & cutoff и point-in-time \\
\texttt{model}, \texttt{provider} & вызов & цена и признаки \\
\texttt{input\_tokens}, \texttt{output\_tokens} & вызов & стоимость по правилам \\
\texttt{reasoning\_tokens} & вызов & бюджет и поведение \\
\texttt{cached\_tokens}, \texttt{cache\_creation\_tokens} & вызов & биллинг и skew \\
\texttt{context\_tokens\_before}, \texttt{context\_growth} & вызов & прогноз TokenCast \\
\texttt{tool\_name}, \texttt{tool\_latency\_ms} & вызов & признаки шага \\
\texttt{status}, \texttt{error\_type} & шаг & исход и цензурирование \\
\texttt{phase} & шаг & активность SDLC \\
\texttt{is\_copied\_context} & шаг & исключение из обучения \\
\texttt{reward}, \texttt{solved} & прогон & метка задачи \\
\texttt{rate\_card\_version} & прогон & пересчёт стоимости \\
\texttt{code\_commit}, \texttt{harness\_version} & прогон & провенанс \\
\end{longtable}
}

\subsection{Что это меняет в решении}

Формат не требует менять алгоритм выбора задач. Он требует изменить
наблюдаемость: вместо итоговых \texttt{cost}/\texttt{spent} харнесс должен
писать поток событий и материализовать таблицу шагов. Это небольшое
изменение поведения, но оно затрагивает \texttt{06\_trajectories.md} и
\texttt{04\_budget\_and\_selection.md}. Спецификацию изменения стоит писать
после эксперимента: сначала собрать 20--50 прогонов, проверить, что
восстановление стоимости из сырых счётчиков совпадает с фактическим
расходом, и обучить простой baseline на таблице шагов. Если baseline
показывает, что шаговых признаков достаточно, формат фиксируется в
спецификации; если нет — набор полей расширяется до эксперимента, а не
после.

\subsection{Следующий шаг}

Эксперимент в \texttt{research/experiments/}: собрать траектории на
доступных задачах, записать их в предложенном формате, пересчитать
стоимость и сравнить три baseline — константа, среднее по задаче и
регрессия на шаговых признаках. Замеры: MAE и покрытие интервалов на
отложенных задачах, доля восстановимых вызовов и токенов, объём хранения,
время сборки Parquet. Отрицательный результат тоже фиксируется: если
формат не даёт прироста над средним по задаче, менять поведение харнесса
ради сбора датасета преждевременно.

\section{Ограничения}

\begin{itemize}
  \item Две базы; Semantic Scholar и Google Scholar недоступны (HTTP 403).
  \item Просматривались верхние записи по релевантности, а не вся выдача;
        для OpenAlex полная выдача по полнотекстовому поиску не
        просматривалась.
  \item Часть источников — препринты, стандарты и технические отчёты, а не
        рецензируемые статьи; они помечены в \texttt{sources.md}.
  \item Данные о форматах извлекались из аннотаций и документации; полные
        тексты читались выборочно.
  \item Downstream-метрики предсказателей стоимости измерены на разных
        наборах (SWE-bench, SpreadsheetBench, SkillsBench, ChatDev) и
        несопоставимы; мета-анализ невозможен.
  \item Нет работы, которая проверяет именно формат хранения на качестве
        предсказателя стоимости; вывод о двухслойной схеме — синтез, а не
        воспроизведённый результат.
  \item Требования соревнования (10 входных = 1 выходной, кеш не
        учитывается) могут отличаться от правил, по которым считались
        стоимости в цитируемых работах; поэтому сырые счётчики важнее
        готового \texttt{cost\_usd}.
\end{itemize}

\bibliographystyle{plain}
\bibliography{refs}

\end{document}
